\section{Implementation}\label{sec:impl}

\subsection{Compute and agentic workflow}\label{sec:workflow}

\paragraph{Compute.} Training and evaluation ran under Slurm on Northeastern's Explorer cluster~\citep{explorerdocs}:
\begin{itemize}[leftmargin=*]
  \item \code{gpu} (H200/A100, 8\,h limit) for training;
  \item \code{sharing} (1\,h) for smoke tests and evaluations;
  \item \code{short} for CPU jobs.
\end{itemize}
Datasets, environments and checkpoints stay on \code{/scratch}; only code, metrics, predictions and figures travel to the laptop, which had about 7\,GB free.

\paragraph{Agents.} Claude Code~\citep{claudecode}, running Claude Opus 5.5, acted as an \emph{orchestrator} (\cref{fig:workflow}). After the student approved a written plan, it launched four parallel workers, each with a self-contained brief:
\begin{itemize}[leftmargin=*]
  \item \emph{repro}: data, environments, released checkpoints and retraining;
  \item \emph{improve}: MMPose experiments;
  \item \emph{analysis}: the evaluator, statistics, figures and generated tables;
  \item \emph{report}: this document and the README.
\end{itemize}
A fresh-context reviewer checks the final state against the rubric.

\paragraph{Live files.} The orchestrator owns three of them:
\begin{itemize}[leftmargin=*]
  \item \code{SKILL.md}: how to operate the cluster;
  \item \code{AGENT.md}: the lab notebook of asks, attempts, failures and dated decisions D1, D2, \dots;
  \item \code{MANIFEST.md}: datasets, checkpoints, configs, and the run behind every number.
\end{itemize}
Workers log commands, job IDs and errors in \code{notebook/<stream>.md}, and register runs in \code{results/registry.csv}.

\paragraph{Shared rules.} All agents load \code{CLAUDE.md}, which requires:
\begin{itemize}[leftmargin=*]
  \item no compute on login nodes;
  \item no selection on Test100;
  \item one evaluator for every number;
  \item secrets only in a git-ignored \code{.env};
  \item timestamps from \code{date}.
\end{itemize}

\paragraph{Cluster loop.} The agent connects with \code{ssh -o BatchMode=yes -o ConnectTimeout=15} and wraps slow login-node commands in \code{timeout}, so failures surface quickly. For each job it:
\begin{enumerate}[leftmargin=*]
  \item pushes code with tar over ssh;
  \item submits an \code{sbatch} script, smoke-tested first on \code{sharing};
  \item polls \code{squeue}, \code{sacct} and the logs;
  \item pulls back only small artifacts.
\end{enumerate}
Each run's directory \code{runs/<RUN_ID>/} holds its resolved config, logs, predictions, evaluator outputs, and a \code{run.json} with the commit and tree SHAs, Slurm job ID, node, GPU, seed and checkpoint SHA256.

\begin{figure}[t]
  \centering
  \resizebox{0.92\linewidth}{!}{% Agentic workflow: human <-> orchestrator <-> worker agents, and the agent <-> Explorer/Slurm loop.
\begin{tikzpicture}[
  font=\sffamily\scriptsize,
  >={Stealth[length=4pt]},
  box/.style={draw, rounded corners=2pt, align=center, inner sep=3pt, line width=0.5pt},
  human/.style={box, fill=realC!18, draw=realC!80!black, minimum width=27mm, minimum height=10mm},
  agent/.style={box, fill=poseC!12, draw=poseC!80!black, minimum width=40mm, minimum height=10mm},
  worker/.style={box, fill=poseC!5, draw=poseC!70!black, minimum width=13.5mm, minimum height=7mm},
  doc/.style={box, fill=white, draw=black!55, minimum width=21mm, font=\ttfamily\tiny, inner sep=2pt},
  clus/.style={box, fill=domC!10, draw=domC!80!black, minimum width=44mm, minimum height=10mm},
  flow/.style={->, line width=0.6pt, draw=black!75},
  lab/.style={font=\sffamily\tiny, text=black!80, align=center},
]
% ---- top row: human -- orchestrator -- cluster login ---------------------------------------
\node[human] (human) at (0,0) {\textbf{student (human)}\\{\tiny sets goals, approves, reviews}};
\node[agent] (orch) at (5.6,0) {\textbf{orchestrator}: Claude Code\\{\tiny (Claude Opus 5.5) on the MacBook}};
\node[clus] (login) at (12.3,0) {\textbf{Explorer login node}\\{\tiny \texttt{ssh -o BatchMode=yes}, \texttt{timeout}; no compute}};
\draw[flow] ([yshift=2mm]human.east) -- node[lab, above] {asks, course guide,\\plan approval, steering} ([yshift=2mm]orch.west);
\draw[flow] ([yshift=-2mm]orch.west) -- node[lab, below] {plan, questions,\\progress, results} ([yshift=-2mm]human.east);
\draw[flow, poseC!80!black] ([yshift=2mm]orch.east) -- node[lab, above] {code: tar over ssh\\jobs: \texttt{sbatch}} ([yshift=2mm]login.west);
\draw[flow, domC!80!black] ([yshift=-2mm]login.west) -- node[lab, below] {\texttt{squeue}, \texttt{sacct}, logs;\\small result files only} ([yshift=-2mm]orch.east);
% ---- second row: workers, Slurm ---------------------------------------------------------------
\node[worker] (w1) at (3.2,-2.1) {repro};
\node[worker] (w2) at (4.8,-2.1) {improve};
\node[worker] (w3) at (6.4,-2.1) {analysis};
\node[worker] (w4) at (8.0,-2.1) {report};
\foreach \w in {w1,w2,w3,w4} \draw[flow] (orch.south) -- (\w.north);
\node[worker, dashed, minimum width=46mm] (rev) at (5.6,-3.25) {fresh-context reviewer (final check)};
\node[clus] (slurm) at (12.3,-2.1) {\textbf{Slurm}\\{\tiny \texttt{gpu}: H200/A100, 8 h \quad \texttt{sharing}: 1 h \quad \texttt{short}: CPU}};
\node[clus] (scratch) at (12.3,-3.85) {\textbf{\texttt{/scratch/.../fidip}}\\{\tiny data, envs, checkpoints, \texttt{runs/<RUN\_ID>/}}};
\draw[flow] (login) -- node[lab, right] {\texttt{sbatch} (smoke test on \texttt{sharing})} (slurm);
\draw[flow] (slurm) -- node[lab, right] {\texttt{run.json}, logs, preds, metrics} (scratch);
% ---- left column: the live files -------------------------------------------------------------
\node[doc] (skill) at (0,-1.55) {SKILL.md};
\node[doc, below=0.8mm of skill] (agentmd) {AGENT.md};
\node[doc, below=0.8mm of agentmd] (manifest) {MANIFEST.md};
\node[doc, below=3.5mm of manifest] (nb) {notebook/<stream>.md};
\node[doc, below=0.8mm of nb] (reg) {results/registry.csv};
\node[lab, anchor=south] at (skill.north) {kept live by the agent};
\draw[flow, black!55] (orch.south west) -- (agentmd.east);
\draw[flow, black!55] (w1.south) |- node[lab, pos=0.75, below] {per-stream detail} (nb.east);
\end{tikzpicture}
}
  \caption{Agentic workflow. The student sets goals, supplies the course material and approves the plan. The orchestrator keeps the three live files current and delegates to parallel workers. All compute goes through Slurm, and only small artifacts return to the laptop.}
  \label{fig:workflow}
\end{figure}

\subsection{Environments and compatibility patches}\label{sec:env}

\paragraph{Environments.} The upstream \code{fidip_env.yml} cannot be built: its pip stage pins \code{grpcio==1.48.2}, which has no Python~3.12 wheel, and \code{mkl-*} packages with no wheels. We built three environments instead:
\begin{itemize}[leftmargin=*]
  \item \emph{FiDIP:} pip-only, Python~3.12, PyTorch~2.3.1 (CUDA~12.1), NumPy~1.26.4, with every imported package pinned to the yml's version.
  \item \emph{MMPose:} Python~3.10, PyTorch~2.1.2, MMEngine~0.10.7, MMCV~2.1.0, MMPretrain~1.2.0 and MMPose~1.3.2~\citep{mmpose2020}, from conda-forge, with \code{numpy<2} and \code{opencv-python<4.12}.
  \item \emph{Analysis:} CPU-only, with the same NumPy and pycocotools pins as FiDIP.
\end{itemize}
Package lists are in \code{envs/} and checkpoint hashes in \code{MANIFEST.md}.

\paragraph{Upstream patches.} \Cref{app:patches} lists six upstream traps and our patches.
The most consequential trap, the README training command, needs no code change: \code{TRAIN.CHECKPOINT} and \code{TRAIN.LAMBDA} are passed on the command line.
The course guide's \code{gpu_test.py}, \code{tools/demo.py} and \code{tools/streamlit_demo.py} (and \code{fidip_simple.yml}) do not exist upstream. We wrote the GPU check and both demos from the guide's documented interface.

\subsection{Data preparation and the split protocol}\label{sec:data}

\paragraph{Download.} A Slurm job (\code{scripts/get_data.sh}):
\begin{enumerate}[leftmargin=*]
  \item downloads both public SyRIP archives, which needed no password;
  \item verifies their SHA256;
  \item extracts them on a compute node, because the login node killed \code{unzip};
  \item checks file counts against the archives;
  \item builds the upstream layout from symlinks and writes the split files.
\end{enumerate}
The adult \code{train_pre_infant} set (2029 real and 2028 synthetic images), described upstream as available on request, ships inside \code{syrip_for_train.zip}, so the classifier pre-stage can be run faithfully.
\code{SyRIP.zip}'s \code{images/train_infant/} also holds 300 unannotated real images, 199 of them byte-identical to validate500 images. We use only images listed in an annotation file.

\paragraph{Training labels.} Two annotation files cover the 1200 training images: the official \code{200R_1000S} file in \code{SyRIP.zip}, and \code{train_infant.json} in \code{syrip_for_train.zip}. They agree on the real images but differ on every synthetic one, because the second file carries MINI-RGBD labels:
\begin{itemize}[leftmargin=*]
  \item its source fields point to MINI-RGBD frames;
  \item eyes and ears are zeroed on most entries;
  \item some sizes and boxes do not fit the renders, with a median keypoint offset of 191\,px from the official labels (\code{notebook/repro.md}).
\end{itemize}
The images in the two archives are identical, except that the second archive's \code{validate_infant/test11.jpg} is a copy of \code{test24.jpg}.
Following the README literally therefore trains on mislabelled synthetic data and scores one test image against the wrong picture: the released checkpoint scores \res{sens.test11.ap}~AP on that archive's test images, against \res{rel.fidip.test.ap}~AP on the correct ones.
We use \code{SyRIP.zip}'s images and \code{200R_1000S} labels everywhere (D17); attempt R020 measures what the README-literal path costs.

\paragraph{Val400 (D15).} validate100 is a subset of validate500 (image ids 0--99, with identical files). \textbf{Val400} is validate500 ids 100--499 minus three images that duplicate Test100 images (two exact copies, one near-duplicate), which leaves 397 images.
Exact copies were found by file hash, and near-duplicates by a flip-aware perceptual hash with visual confirmation (\code{scripts/near_dups.py}).

\paragraph{Train/test overlap (D16).} Test100 images \code{test20}, \code{test42}, \code{test46} and \code{test50} are byte-identical to training images 73, 315, 272 and 340, and \code{test7} is a near-duplicate of training image 47.
We keep Test100 as shipped, for comparability with the published numbers (which carry the same overlap). As a sensitivity check we also report the 95-image \emph{Test100-clean} subset, where the released model's AP changes by \res{sens.clean.delta}.
Two further properties of the shipped GT need care.

\paragraph{Placeholder areas (D8).} Every validate500 annotation carries the same \code{area}, 2443298.48 ($\sqrt{\text{area}}\approx1563$\,px), and the training file carries another constant, 14232.77; validate100 instead uses $\text{area}=w\cdot h$ of its box.
Since OKS divides squared errors by the area (\cref{eq:oks}), the shipped Val400 makes almost any prediction look perfect: the released model scores \res{sens.area.const.ap}~AP on it.
We rebuilt Val400, validate500 and the framework-neutral training file with $\text{area}=w\cdot h$ from their own boxes, and left Test100 exactly as shipped. \code{scripts/make_splits.py} asserts this rule, and the evaluator refuses GT files whose areas are all identical.

\paragraph{Different boxes.} On the 100 shared images, validate100's boxes are a median \res{data.boxratio.pct}\% larger than validate500's, while the keypoints agree on 98 images; the other two were re-annotated. Val400 boxes are therefore tighter than Test100's, which means a smaller OKS scale and tighter crops (\cref{sec:thresholds}).

\paragraph{Protocol (D1).} \Cref{fig:protocol} summarises the protocol:
\begin{itemize}[leftmargin=*]
  \item Every choice of epoch, hyperparameter, backbone or recipe is made on Val400.
  \item A configuration reaches Test100 only after a commit lists it, with its Val400 evidence and checkpoint rule, in \code{results/PREREGISTRATION.md}.
  \item The final-evaluation script refuses unregistered or repeated runs, and checks the checkpoint SHA256 against \code{run.json}.
  \item The released checkpoints and the official-recipe retrains are reproduction targets, so they are evaluated on Test100 directly; the retrains' epochs are still chosen on Val400.
\end{itemize}

\paragraph{Why Val400.} It keeps the training data identical to the paper's, and it is almost twice as large as the real training set; cross-validation on the 200 real images could not match that without shrinking training.
It has two costs: Test500 is no longer clean, so we never report it, and Val400 is not the complex-pose subset, so it may rank models differently from Test100 (\cref{sec:thresholds}).

\begin{figure}[t]
  \centering
  \resizebox{0.92\linewidth}{!}{% Split and selection protocol: train on 200R+1000S, select on Val400, report once on Test100.
\begin{tikzpicture}[
  font=\sffamily\scriptsize,
  >={Stealth[length=4pt]},
  seg/.style={draw, line width=0.5pt, minimum height=8mm, align=center, inner sep=2pt},
  pstep/.style={draw, rounded corners=2pt, align=center, inner sep=3pt, line width=0.5pt, minimum height=11mm, fill=white},
  flow/.style={->, line width=0.6pt, draw=black!75},
  lab/.style={font=\sffamily\tiny, text=black!75, align=center},
]
% ---- the data bar (widths proportional to image counts: 1200 | 400 | 100) ----------------------
\def\u{0.0068}  % cm per image
\node[seg, fill=realC!25, draw=realC!80!black, minimum width=200*\u cm, anchor=west] (tr) at (0,0) {200\\real};
\node[seg, fill=synC!20, draw=synC!80!black, minimum width=1000*\u cm, anchor=west] (ts) at (tr.east) {1000 synthetic (SMIL renders)};
\node[seg, fill=poseC!15, draw=poseC!80!black, minimum width=397*\u cm, anchor=west, xshift=3mm] (val) at (ts.east)
  {\textbf{Val400}\\{\tiny 397 images: ids 100--499 minus}\\[-2pt]{\tiny 3 duplicates of Test100}};
\node[seg, fill=lossC!15, draw=lossC!85!black, minimum width=100*\u cm+6mm, anchor=west] (test) at (val.east)
  {\textbf{Test100}\\{\tiny ids 0--99}};
\draw[decorate, decoration={brace, amplitude=4pt}, line width=0.5pt] (tr.north west) -- node[lab, above=4pt] {SyRIP train (official \texttt{200R\_1000S} annotations)} (ts.north east);
\draw[decorate, decoration={brace, amplitude=4pt}, line width=0.5pt] (val.north west) -- node[lab, above=4pt] {validate500 = Test500 (500 real); validate100 $\subset$ validate500} (test.north east);
\node[lab, right=1mm of test.east, anchor=west, text=lossC!85!black, align=left] {incl.\ 5 copies\\of training\\images (D16)};
% ---- the procedure ----------------------------------------------------------------------------
\node[pstep, below=9mm of ts.south west, anchor=north west, xshift=-14mm] (s1)
  {\textbf{1 train}\\candidate configs\\$\times$ seeds};
\node[pstep, right=6mm of s1] (s2) {\textbf{2 select on Val400}\\epoch, hyper-params,\\backbone, recipe};
\node[pstep, right=6mm of s2] (s3) {\textbf{3 pre-register}\\config + ckpt rule\\committed in\\\texttt{PREREGISTRATION.md}};
\node[pstep, right=6mm of s3] (s4) {\textbf{4 one Test100 eval}\\script refuses\\unregistered or\\repeated runs};
\node[pstep, right=6mm of s4, fill=black!4] (s5) {\textbf{5 report}\\\texttt{eval\_coco.py}\\+ paired bootstrap\\vs.\ released FiDIP};
\draw[flow] (s1) -- (s2);
\draw[flow] (s2) -- (s3);
\draw[flow] (s3) -- (s4);
\draw[flow] (s4) -- (s5);
\draw[flow, black!55] ($(tr.south west)!0.5!(ts.south east)$) -- ++(0,-3mm) -| (s1.north);
\draw[flow, poseC!80!black] (val.south) -- ++(0,-3mm) -| ([xshift=4mm]s2.north);
\draw[flow, lossC!85!black] (test.south) -- ++(0,-3mm) -| (s4.north);
% ---- the forbidden edge ----------------------------------------------------------------------------
\draw[->, line width=0.6pt, lossC!85!black, densely dashed] (s4.south) -- ++(0,-5mm) -| node[pos=0.25, below, lab, text=lossC!85!black]
  {no Test100 feedback into steps 1--3 (CLAUDE.md rule; \texttt{final\_test\_eval} checks the registry)} (s1.south);
\node[lab, anchor=north west, align=left] at ([yshift=-11mm]s1.south west)
  {Exception (stated in \texttt{PREREGISTRATION.md}): the released checkpoints and the official-recipe retrains are reproduction targets, not tuned candidates,\\so they are evaluated on Test100 directly. Their epoch is still chosen on Val400.};
\end{tikzpicture}
}
  \caption{Split and selection protocol (bar widths proportional to image counts). Test100 is read once per pre-registered configuration, and nothing learned from it flows back into training or selection.}
  \label{fig:protocol}
\end{figure}

\subsection{One evaluator for every number}\label{sec:eval}

\paragraph{Pipeline.} Every metric comes from \code{eval/eval_coco.py} (v1.0) applied to a prediction JSON, and both frameworks take the same path:
\begin{enumerate}[leftmargin=*]
  \item a loader that accepts twelve prediction encodings;
  \item one rescoring rule, score $=\mathrm{mean}(c_k\,|\,c_k>0.2)$, which is upstream FiDIP's rule and MMPose's \code{bbox_keypoint} mode with a box score of 1 for GT boxes. It keeps the image ranking, and therefore AP, comparable;
  \item a verbatim port of upstream OKS-NMS, which is a no-op with one box per image;
  \item pycocotools 2.0.8 \code{COCOeval}~\citep{cocoapi}, with the COCO sigmas and the GT area.
\end{enumerate}
Besides the usual metrics it writes per-threshold AP, \APstrict{} and per-image OKS. It also writes per-keypoint errors normalised by $\sqrt{wh}$ of the GT box (the OKS scale on Test100), with PCK@\{0.05, 0.1, 0.2\}.

\paragraph{The annotation-id-0 artifact (D10).} In Test100, annotation ids equal image ids, so \code{test0.jpg} carries id~0. \code{COCOeval.evaluateImg} stores the matched GT id in \code{dtMatches}, and \code{accumulate} treats a stored 0 as ``unmatched''. The infant in \code{test0.jpg} can therefore never be a true positive, and its detection always counts as a false positive.
This caps AR at 0.99, and even a perfect model's AP at every threshold is
\[
\frac{r+0.99\,(100-r)}{101}\in[0.980,\,0.990],
\]
where $r$ is the rank of test0's score (unit-tested).
The published AP$_{50}$ = AP$_{75}$ = 98.5 falls inside this band, so on the 99 matchable images the published models are already perfect at loose thresholds. The near-constant AR$_{50}$ = 99.0 of the upstream table has the same cause.

The size of this penalty depends only on where test0's score ranks: between \res{sens.id0.min} and \res{sens.id0.max}~AP at realistic accuracy.
Our headline \texttt{paper} protocol keeps the ids as shipped, so it stays comparable with 93.6. A \texttt{fixed} protocol shifts the ids off 0. An improvement claim counts only if its paired CI excludes 0 under \emph{both} protocols, and every model's test0 rank is reported (D13). On Val400, which has no id~0, the two protocols coincide.
This is a property of the evaluation tooling, not an error in the published work: anyone scoring this file with pycocotools gets the same cap.

\paragraph{MMPose's evaluator.} MMPose's \code{CocoMetric} uses xtcocotools~\citep{xtcocotools}, which marks unmatched detections with $-1$. Its printed Test100 AP therefore equals our \texttt{fixed} protocol, \res{sens.xtcoco.delta}~AP away from \texttt{paper} on identical predictions.
Comparing it with 93.6 would overstate any gain, which is why only \code{eval_coco.py} numbers are compared (D2).

\paragraph{Verification and statistics.}
\begin{itemize}[leftmargin=*]
  \item \emph{Unit tests.} 24 tests cover the GT-as-prediction case, the id-0 ceiling, noise monotonicity, encoding invariance, rescoring, OKS against \code{COCOeval} internals, and NMS against upstream (\cref{app:tests}).
  \item \emph{Cross-check.} On the released checkpoint's predictions, upstream's own test script reports \res{rel.fidip.test.ap.upstream}~AP and ours \res{rel.fidip.test.ap}~AP.
  \item \emph{Uncertainty.} \code{eval/bootstrap_ci.py} resamples the 100 test images ($B=\res{boot.n}$), recomputes AP exactly, and reports paired 95\% percentile CIs of seed-averaged AP differences and Holm-adjusted one-sided $p$-values~\citep{efron1993bootstrap}.
\end{itemize}
